\section{A research agenda}\label{sec:agenda}

The SMoE is minimal by design. Each extension below keeps the white-box property (closed-form
predictive functionals, exact attribution and likelihood-based inference) and each is a
self-contained research question.

\paragraph{1. Multivariate SMoE for allocation.}
Replace the univariate experts by $\N_d(\bm B_k\tilde\bx,\bm\Sigma_k(\bv))$ with regime-specific
covariance matrices, e.g.\ $\bm\Sigma_k=\bm D_k\bm R_k\bm D_k$ with log-linear $\bm D_k$. The
predictive covariance decomposes as $\sum_k\pi_k\bm\Sigma_k+\sum_k\pi_k(\bm\mu_k-\bar{\bm\mu})(\bm\mu_k-\bar{\bm\mu})^\top$,
and the second term is a rank-$(K-1)$ ``regime-disagreement'' factor structure. Open questions:
identification with many assets, shrinkage of $\bm R_k$, and whether regime-disagreement factors are
priced.

\paragraph{2. Decision-focused estimation.}
\Cref{sec:lagging} traces the lagging paradox to a mismatch between the training loss and the use of the
forecast. A natural next step estimates $\btheta$ by maximising realised utility,
$\max_{\btheta}\sum_tU(w_t(\btheta)y_t)$ with $w_t$ given by \cref{prop:band}, or a convex combination of
utility and likelihood. The questions are statistical (consistency of a non-likelihood M-estimator
whose criterion is non-smooth because of the no-trade band) and economic (does it beat
predict-then-optimise out of sample?).

\paragraph{3. Language-model gates.}
Gate covariates need not be prices. Policy statements, regulatory announcements and news can be
mapped by a large language model to a small set of scores (e.g.\ policy stance, liquidity support,
risk-off intensity) that enter $\bz_t$. The SMoE then answers a sharp question (does text move
regime probabilities beyond what prices already imply?) with the score test of \cref{prop:lm-gate},
which needs only the price-based fit and can screen many text scores at once, followed by a Wald test
on the selected $\bm\gamma$ coefficients. The fidelity ratio then measures how much of an LLM-based
black-box forecaster the text-augmented SMoE explains. Because language models also carry look-ahead
in their pretraining, the scores should come from models with dated corpora (item 8).

\paragraph{4. Online Bayesian SMoE.}
Quarterly re-estimation is a crude response to non-stationarity. Sequential Monte Carlo or online
variational EM would update $\btheta$ daily with forgetting, give posterior uncertainty about the
regime parameters themselves, and allow the gate to reflect parameter uncertainty in the predictive
distribution. The adaptive conformal layer of \cref{thm:aci} would then calibrate a Bayesian rather
than a frequentist predictive.

\paragraph{5. Intraday SMoE.}
With 5-minute data the same model can forecast intraday returns and realised variance, with experts
for opening auction, lunch break and closing auction dynamics specific to the Chinese market, and gates
driven by order-flow imbalance. Microstructure noise makes the distributional target essential.

\paragraph{6. Realised-variance experts.}
\Cref{sec:rv} found that a single-regime Realized GARCH beats every SMoE variant in log score once
5-minute data are available, and that part of the mixture's advantage with range-based covariates is
insurance against a noisy variance measure. The obvious synthesis gives each expert its own
measurement equation and variance recursion,
$\log\sigma^2_{kt}=\omega_k+\beta_k\log\sigma^2_{k,t-1}+\gamma_k\log\mathrm{RV}_{t-1}$, with the
gate deciding which recursion governs tomorrow. The CM-step for the variance block is then no longer
a static Gamma regression but a filtering problem, and identification of regime-specific persistence
from a single realised series is open. Extending realised measures to the other eighteen series
requires intraday data we do not have.

\paragraph{7. Yardstick for foundation models.}
As time-series foundation models grow, the fidelity ratio becomes a general tool: for each model and
asset it reports how much of the model's out-of-sample density-forecast gain is reproducible by a
three-regime, ten-covariate white box. A benchmark of $\Phi$ across model families, horizons and
asset classes would make ``how black is the black box?'' a measurable property. Newer models with
covariate and multivariate inputs, such as TimesFM~3.0, can be given the SMoE's own covariates, which
turns the yardstick into a direct comparison of what the network and the white box do with the same
information.

\paragraph{8. Point-in-time foundation models.}
The release-date test of \cref{sec:lockbox} can only bound pretraining look-ahead from one side,
because a model's release date is an upper bound on its data cut-off. Foundation models pretrained on
corpora that end at known, staggered dates would turn the test into an exact event study: the same
architecture trained to vintages $C_1<C_2<\dots$ and evaluated on the same returns identifies how much
of an apparent forecasting gain is memory. \citet{sarkar2024lookahead} make the same argument for
language models.

